%-------------------------------------------------------------------------------
% File:		Homework.tex
% Author:	Igor Janjic, Brian Hilnbrand, Danny Duangphachanh, Leah
%		Krynitsky  
% Description:	[ECE 4534] Embedded Systems Design
%		Background Research Assignment
%%------------------------------------------------------------------------------ 

%\documentclass{Book}

%-------------------------------------------------------------------------------
% File:		Preamble.tex
% Author:	Igor Janjic
% Description:	Defines which packages to use.
%%------------------------------------------------------------------------------

\documentclass[11pt]{article}

% Encodes font as T1.
\usepackage[T1]{fontenc}

% Pretty much all of the ams maths packages.
\usepackage{amsmath,amsthm,amssymb,amsfonts}

% Allows you to manipulate the page a bit.
\usepackage[margin=0.75in]{geometry}

% Removes paragraph indentation.
\usepackage{parskip}

% Allows inclusion of graphics easily.
\usepackage{graphicx}

% Provides ways to make nice looking tables.
\usepackage{booktabs}

% Allows rotation of tables and figures.
\usepackage{rotating}

% Allows color.
\usepackage{color}

% Allows shading of table cells.
\usepackage{colortbl}

% Allows hyphenatable letterspacing, underlining, and highlighting.
\usepackage{soul}

% Allows extra text symbols.
\usepackage{textcomp}

% Allows manipulation of floats.
\usepackage{float}

% Provides commands to make subfigures.
\usepackage{subfigure}

% Typesets URLs sensibly with tt font, clickability in PDFs, and not breaking across lines.
\usepackage{url}

% Makes references hyperlinks in PDF output.
\usepackage{hyperref}

% Provides good access to colors.
\usepackage{color}

% Provides good access to colors.
\usepackage{xcolor}

% Provides for programming code to be inserted into the document.
\usepackage{listings}

% Provides support for table coloring.
\usepackage{xcolor}

% Provides vector graphics functionality.
\usepackage{tikz}
\usetikzlibrary{shapes.geometric, arrows, positioning}

% Provides circuit vector graphics functionality.
\usepackage{circuitikz}

% Provides caption manipulation.
\usepackage{caption}

% Makes sure tables don't float about if they have the designation [H].
\restylefloat{table}

% Defines the color gray.
\definecolor{gray}{rgb}{0.5, 0.5, 0.5}


%-------------------------------------------------------------------------------
% File:		Definitions.tex
% Author:	Igor Janjic
% Description:	Contains various commands.
%%------------------------------------------------------------------------------

%-------------------------------------------------------------------------------
% General Definitions

% Defines a command for skipping a line.
\def\wl{\par\vspace{\baselineskip}}

% Defines a command for hiding text.
\newcommand{\hide}[1]{}

% Defines a command for a shortened commonly used phrases.
\newcommand{\ie}{\text{i.e.}}
\newcommand{\eg}{\text{e.g.}}
\newcommand{\ex}{\text{e.x.}}
\newcommand{\ans}{\textbf{Answer: }}
\newcommand{\pro}{\textbf{Proof: }}

% Defines a command to mark to-do items in red.
\newcommand{\todo}[1] {\textbf{\textcolor{red}{#1}}}

% Defines commands for calligraphic font.
\newcommand{\calA}{\mathcal{A}}
\newcommand{\calB}{\mathcal{B}}
\newcommand{\calC}{\mathcal{C}}
\newcommand{\calD}{\mathcal{D}}
\newcommand{\calE}{\mathcal{E}}
\newcommand{\calF}{\mathcal{F}}
\newcommand{\calG}{\mathcal{G}}
\newcommand{\calH}{\mathcal{H}}
\newcommand{\calI}{\mathcal{I}}
\newcommand{\calJ}{\mathcal{J}}
\newcommand{\calK}{\mathcal{K}}
\newcommand{\calL}{\mathcal{L}}
\newcommand{\calM}{\mathcal{M}}
\newcommand{\calN}{\mathcal{N}}
\newcommand{\calO}{\mathcal{O}}
\newcommand{\calP}{\mathcal{P}}
\newcommand{\calQ}{\mathcal{Q}}
\newcommand{\calR}{\mathcal{R}}
\newcommand{\calS}{\mathcal{S}}
\newcommand{\calT}{\mathcal{T}}
\newcommand{\calU}{\mathcal{U}}
\newcommand{\calV}{\mathcal{V}}
\newcommand{\calW}{\mathcal{W}}
\newcommand{\calX}{\mathcal{X}}
\newcommand{\calY}{\mathcal{Y}}
\newcommand{\calZ}{\mathcal{Z}}

% Defines commands for bold font.
\newcommand{\bbA}{\mathbb{A}}
\newcommand{\bbB}{\mathbb{B}}
\newcommand{\bbC}{\mathbb{C}}
\newcommand{\bbD}{\mathbb{D}}
\newcommand{\bbE}{\mathbb{E}}
\newcommand{\bbF}{\mathbb{F}}
\newcommand{\bbG}{\mathbb{G}}
\newcommand{\bbH}{\mathbb{H}}
\newcommand{\bbI}{\mathbb{I}}
\newcommand{\bbJ}{\mathbb{J}}
\newcommand{\bbK}{\mathbb{K}}
\newcommand{\bbL}{\mathbb{L}}
\newcommand{\bbM}{\mathbb{M}}
\newcommand{\bbN}{\mathbb{N}}
\newcommand{\bbO}{\mathbb{O}}
\newcommand{\bbP}{\mathbb{P}}
\newcommand{\bbQ}{\mathbb{Q}}
\newcommand{\bbR}{\mathbb{R}}
\newcommand{\bbS}{\mathbb{S}}
\newcommand{\bbT}{\mathbb{T}}
\newcommand{\bbU}{\mathbb{U}}
\newcommand{\bbV}{\mathbb{V}}
\newcommand{\bbW}{\mathbb{W}}
\newcommand{\bbX}{\mathbb{X}}
\newcommand{\bbY}{\mathbb{Y}}
\newcommand{\bbZ}{\mathbb{Z}}

% Defines commands for bold face font.
\newcommand{\bfA}{\mathbf{A}}
\newcommand{\bfB}{\mathbf{B}}
\newcommand{\bfC}{\mathbf{C}}
\newcommand{\bfD}{\mathbf{D}}
\newcommand{\bfE}{\mathbf{E}}
\newcommand{\bfF}{\mathbf{F}}
\newcommand{\bfG}{\mathbf{G}}
\newcommand{\bfH}{\mathbf{H}}
\newcommand{\bfI}{\mathbf{I}}
\newcommand{\bfJ}{\mathbf{J}}
\newcommand{\bfK}{\mathbf{K}}
\newcommand{\bfL}{\mathbf{L}}
\newcommand{\bfM}{\mathbf{M}}
\newcommand{\bfN}{\mathbf{N}}
\newcommand{\bfO}{\mathbf{O}}
\newcommand{\bfP}{\mathbf{P}}
\newcommand{\bfQ}{\mathbf{Q}}
\newcommand{\bfR}{\mathbf{R}}
\newcommand{\bfS}{\mathbf{S}}
\newcommand{\bfT}{\mathbf{T}}
\newcommand{\bfU}{\mathbf{U}}
\newcommand{\bfV}{\mathbf{V}}
\newcommand{\bfW}{\mathbf{W}}
\newcommand{\bfX}{\mathbf{X}}
\newcommand{\bfY}{\mathbf{Y}}
\newcommand{\bfZ}{\mathbf{Z}}

% Defines commands for lists.
\newcommand{\be}{\begin{enumerate}}
\newcommand{\ee}{\end{enumerate}}
\newcommand{\ba}{\begin{align}}
\newcommand{\ea}{\end{align}}
\newcommand{\bas}{\begin{align*}}
\newcommand{\eassss}{\end{align*}}

%-------------------------------------------------------------------------------
% Math Definitions

% Defines a command for a function.
\newcommand{\func}[2]{#1 \left( #2 \right)}

% Defines a command for an integral.
\newcommand{\dint}[4]{\int_{#1}^{#2}#3\;d#4}

% Defines commands for objects.
\newcommand{\csum}[2]{\sum_{#1}^{#2}}
\newcommand{\cint}[2]{\int_{#1}{#2}}
\newcommand{\cprod}[2]{\prod_{#1}^{#2}}

% Defines commands for various parentheses.
\newcommand{\paren}[1]{\left( #1 \right)}
\newcommand{\sqprn}[1]{\left[ #1 \right]}
\newcommand{\tlprn}[1]{\left\{ #1 \right\}}

% Defines a command for the absolute value of an expression.
\newcommand{\abs}[1]{\left| #1 \right|}

% Defines commands for various angled expressions.
\newcommand{\tr}[2]{\left<\left< #1 , #2 \right>\right>}
\newcommand{\trs}[2]{\left< #1 , #2 \right>}
\newcommand{\iprod}[2]{\left\langle #1 , #2 \right\rangle}

% Defines commands for various deltas.
\newcommand{\del}{\partial}
\newcommand{\dd}[2]{\frac{d#1}{d#2}}
\newcommand{\ddel}[2]{\frac{\del#1}{\del#2}}
\newcommand{\ddell}[2]{\frac{\del^2#1}{{\del#2}^2}}
\newcommand{\ddelll}[3]{\frac{\del^2#1}{{\del#2}{\del#3}}}

% Defines a command for evaluating an expression at a particular value.
\newcommand{\eval}[3]{\left.#1\right|_{#2}^{#3}}

% Defines commands for flooring and ceiling.
\newcommand{\ceil}[1]{\ensuremath{\left\lceil#1\right\rceil}}
\newcommand{\floor}[1]{\ensuremath{\left\lfloor#1\right\rfloor}}

% Defines commands for argmin and argmax.
\newcommand{\argmin}{\mathop{\mathrm{argmin}}}
\newcommand{\argmax}{\mathop{\mathrm{argmax}}}

% Defines a command for mapping.
\newcommand{\map}{\mathop{\mathrm{map}}}

% Defines a command for an inverse.
\newcommand{\inv}{^{-1}}

% Defines a command for the partial derivative operator.
\newcommand{\pd}[2]{\frac{\partial #1}{\partial #2}}

% Defines commands for function names.
\newcommand{\zeros}{\mathrm{zeros}}
\newcommand{\rank}{\mathrm{rank}}
\newcommand{\TR}{\mathrm{tr}}
\newcommand{\vol}{\mathrm{vol}}
\newcommand{\sgn}{\mathrm{sgn}}
\newcommand{\acos}{\mathrm{acos}}
\newcommand{\intr}{\mathrm{int}\,}
\newcommand{\extr}{\mathrm{ext}\,}
\newcommand{\dom}{\mathrm{dom}\,}
\newcommand{\conv}{\mathrm{conv}\hspace{1pt}}
\newcommand{\cone}{\mathrm{cone}\hspace{1pt}}

% Defines commands for symbols in probability and statistics.
\newcommand{\convd}{\overset{{\scriptsize d}}{\to}}
\newcommand{\convp}{\overset{{\scriptsize p}}{\to}}
\newcommand{\unif}{\mathrm{Uniform}}
\newcommand{\bern}{\mathrm{Bernoulli}}
\newcommand{\negbin}{\mathrm{NegBinom}}
\newcommand{\Binom}{\mathrm{Bin}}
\newcommand{\pois}{\mathrm{Poisson}}
\newcommand{\ndistr}{\mathrm{n}}
\newcommand{\Bdistr}{\mathrm{Beta}}
\newcommand{\Gdistr}{\mathrm{Gamma}}
\newcommand{\iGdistr}{\mathrm{InvGamma}}
\newcommand{\MSE}{\mathrm{MSE}}
\newcommand{\MAE}{\mathrm{MAE}}
\newcommand{\ARE}{\mathrm{ARE}}
\newcommand{\se}{\mathrm{se}}
\newcommand{\mexp}{\text{{\bf E}}}
\newcommand{\var}{\mathrm{Var}}
\newcommand{\cov}{\mathrm{Cov}}
\newcommand{\corr}{\mathrm{Cor}}
\newcommand{\bW}{\bs{W}}
\newcommand{\bw}{\bs{w}}
\newcommand{\bX}{\bs{X}}
\newcommand{\bx}{\bs{x}}
\newcommand{\bY}{\bs{Y}}
\newcommand{\by}{\bs{y}}
\newcommand{\bZ}{\bs{Z}}
\newcommand{\bz}{\bs{z}}
\newcommand{\Xbar}{\overline{X}}
\newcommand{\ev}{\mathrm{Ev}}
\newcommand{\pr}{\boldsymbol{\mathrm{P}}}


%-------------------------------------------------------------------------------
% File:		Programming.tex
% Author:	Igor Janjic
% Description:	Defines programming environments.
%%------------------------------------------------------------------------------

\lstnewenvironment{matlab}{
\lstset{
    	language = matlab,	
	basicstyle = \footnotesize\ttfamily,
	numbers = left,
	numberstyle = \tiny\color{gray},
	stepnumber = 1,
	numbersep = 5pt,
 	backgroundcolor = \color{white},
 	showspaces = false,
  	showstringspaces = false,
  	showtabs = false,
  	frame = single,
  	rulecolor = \color{black},
  	tabsize = 4,
  	breaklines = true,
  	breakatwhitespace = false,
  	title = \lstname,
    	upquote = true,
    	aboveskip = {1.5\baselineskip},
    	columns=fixed,
   	extendedchars = true,
   	prebreak = \raisebox{0ex}[0ex][0ex]{\ensuremath{\hookleftarrow}},
    	identifierstyle=\ttfamily,
    	keywordstyle=\color[rgb]{0,0,1},
    	commentstyle=\color[rgb]{0.133,0.545,0.133},
    	stringstyle=\color[rgb]{0.627,0.126,0.941}
}}{}

\lstnewenvironment{tex}{
\lstset{
    	language = TeX,	
	basicstyle = \footnotesize\ttfamily,
	numbers = left,
	numberstyle = \tiny\color{gray},
	stepnumber = 1,
	numbersep = 5pt,
 	backgroundcolor = \color{white},
 	showspaces = false,
  	showstringspaces = false,
  	showtabs = false,
  	frame = shadowbox,
  	rulecolor = \color{black},
  	tabsize = 4,
  	breaklines = true,
  	breakatwhitespace = false,
  	title = \lstname,
    	upquote = true,
    	aboveskip = {1.5\baselineskip},
    	columns=fixed,
   	extendedchars = true,
    	identifierstyle=\ttfamily,
	xleftmargin=0.5cm,
	xrightmargin=0.5cm,
    	keywordstyle=\color[rgb]{0,0,1},
    	commentstyle=\color[rgb]{0.133,0.545,0.133},
    	stringstyle=\color[rgb]{0.627,0.126,0.941}
}}{}

\lstnewenvironment{pseudo}{
\lstset{
	basicstyle = \footnotesize,
	numbers = left,
	numberstyle = \tiny\color{gray},
	stepnumber = 1,
	numbersep = 5pt,
 	backgroundcolor = \color{white},
 	showspaces = false,
  	showstringspaces = false,
  	showtabs = false,
  	frame = single,
  	rulecolor = \color{black},
  	tabsize = 4,
  	captionpos = b,
  	breaklines = true,
  	breakatwhitespace = false,
  	title = \lstname,
    	upquote = true,
    	aboveskip = {1.5\baselineskip},
    	columns=fixed,
   	extendedchars = true,
   	prebreak = \raisebox{0ex}[0ex][0ex]{\ensuremath{\hookleftarrow}},
}}{}



\begin{document}

%-------------------------------------------------------------------------------
% File:		    Title.tex
% Author:	    Igor Janjic
% Description:	[ECE 4564] Network Applications Design
%		        Project proposal.
%%------------------------------------------------------------------------------

\begin{titlepage}

\centering
\vspace*{\baselineskip}

\rule{\textwidth}{1.6pt}\vspace*{-\baselineskip}\vspace*{2pt}
\rule{\textwidth}{0.4pt}\\[\baselineskip]

{\LARGE Final Report\\[0.3\baselineskip]}

\rule{\textwidth}{0.4pt}\vspace*{-\baselineskip}\vspace{3.2pt}
\rule{\textwidth}{1.6pt}\\[\baselineskip]

\wl

\scshape Due December 12, 2014 at 11:55 PM.
{\small 
\\[\baselineskip]\par}

\vfill

Created by:\\[0.2\baselineskip]
{Brian Hilnbrand:     \texttt{brhiln@vt.edu}}\\[0.2\baselineskip]
{Danny Duangphachanh: \texttt{bboydd@vt.edu}}\\[0.2\baselineskip]
{Igor Janjic:         \texttt{ijanjic@vt.edu}}\\[0.2\baselineskip]
{Leah Krynitsky:      \texttt{leah8@vt.edu}}\\[0.4\baselineskip]
{\small \today}\\[0.8\baselineskip]
{\small [ECE 4534] Embedded Systems Design}\\[0.2\baselineskip]
{\small\itshape Virginia Polytechnic Institute and\\ State University}\\[0.2\baselineskip]

\begin{center}
	\includegraphics[scale=0.35]{Images/Logo}
\end{center}

\end{titlepage}


\subsection*{Background}
Of the first autonomous robots, none are more famous than Grey Walter's tortoises Elmer and Elsie. These two robots were very simple analog circuits equipped with three wheels for locomotion and light sensors which they used to navigate their environment and move toward illuminated recharging stations when low on power \cite{Walter:1950}. Since then, the design of autonomous robots has shifted toward the processing of sensory data using algorithms based in terms of digital computation. The remaining paper focuses on this last view of autonomous robotics.

\subsection*{Robot Framework Considerations}
The three primitives of robotics: plan, sense, act, as well as their relationships imply a robotic framework. Since the design specifications and hardware limitations are given, the robotics framework must then be determined by judging each framework's strengths and weaknesses as well as the compatibility between the framework and the choice of available sensors (described later). Thus, we are purposefully limiting our discussion to only those frameworks that coincide with usability of the various sensors available.

The framework must be able to handle quickly processing large amounts of complex information. In the 1960’s, development of a roving vehicle for surface exploration of the planet Mars demonstrated the need to have the framework on the vehicle to help eliminate lag. This implementation proved useful when tested with terrain that was not smooth and level. \cite{Hogle}

Additionally, the framework must be able to handle the prioritization of tasks. Since the robot has multiple goals that it is trying to achieve, often at the same time and often contradictory to each other, the control system framework must be able to rank goals by priority and focus computational power to solving high-priority goals first while at the same time reserving enough computational power to service low-priority goals.

The more sensors the robot has, the more information it is able to gather and then process to help it achieve its goals. The framework must be able to handle error in the sensors, overlapping sensors, inconsistent or failed sensors, and the addition of more sensors, 

\subsection*{Layered Control Paradigm \cite{Brooks:1985}}
This paradigm was designed by Brooks, R. at the MIT Artificial Intelligence Laboratory, which allowed for the design of various levels of intelligence in autonomous robots by building the control system in layers. Each layer is composed of simple computational modules communicating over low bandwidth channels asynchronously. The problem formulated in the language of control theory is as follows: a group of sensors receive information about the environment which the control system processes, turns into commands, and then sends to the motor controller which moves the robot accordingly.

\subsection*{Choice of Sensors}
The sensors implemented in this design will be used mostly for obstacle detection and navigation.  Some of the most commonly used sensors for these tasks include distance sensors and proximity sensors.  These types of sensors are used to determine the relative distance between the sensor and objects in the environment and can be either passive or active. 

The main type of passive sensor used in object detection is vision based, like a visible spectrum camera.  The images captured by this type of system provide a large amount of information and are easy to interpret. Since passive sensors do not emit any signals, the environmental interference issues that can arise with active sensors are not present.  Although passive sensors are advantageous in this regard, they perform poorly in low light environments since they must rely on natural energy in the environment.  Alternatively, active sensors provide their own energy source for object detection and dramatically outperform passive systems in environments lacking illumination.  Radar, ladar, and sonar are the three main types of active sensors that can be implemented in an object detection system, and all implement a system that emits a signal and interprets reflections to determine information about an object.  \cite{Discant}

Infrared sensors, available as passive or active sensors, are very widely used and offer many advantages over other types of sensors when performing navigation, object avoidance, and line following.  Infrared sensors are also cheap in comparison to other sensor choices and are able to operate in real-time, making them a popular choice for designs like this one.\cite{Calin}

When designing a sensor system, more than one type of sensor can be used when multiple sensors are to be implemented.  As a result, many designs optimize performance by creating a system which includes both active and passive sensors, thus acheiving the advantages of each and eliminating most of the disadvantages that come from using only one type of sensor.  A design which implements a mix of active and passive sensors improves both system performance and object recognition. \cite{Discant}  

\subsection*{Comparison of Algorithms}


\subsection*{Communication Protocols}
In general terms, a message intended for communication between two devices (the protocol) is comprised of two main parts: the header and the payload. The header consists of information about the message, including source or destination id, the nature of the message, and the size of the payload. The payload is the data which is being sent from one device to another. This payload can be formatted in any way, as long as it is consistent in both the source and destination.

Google Protocol Buffers \cite{Google:2014} are used in the automotive industry to format and serialize both small and large packets of information, preparing the payload for transmission to any destination. Once recieved at the destination, the payload can be deserialized and treated as an object, pulling specific pieces of data out of the payload as necessary. A single .proto file is written to define any number of different payloads to send, which is compiled by the protocol buffer compiler to produce libraries to be used for serialization. This provides an easy way to create payloads with lots of information if necessary, which can be created and understood seemlessly by both the source device and the destination device.

In addition, a modified transmission protocol \cite{NSI:2014} should be used to define the order and confirmation of transmission. If the link between two devices is thought of as an open stream of communication, it is wise to set up a series of handshakes and data packets to ensure correct communication. For example, one may want to send the header first to alert the destination device of an incomming payload, what information it will contain, and how long the stream of data is. Once the destination device has processed this header and is ready to recieve information from the source, such as a sensor data point, it will send a handshake, or similar confirmation. Since the destination device already knows the length of the incomming payload, there is no need to analyze it as it is comming in to catch a delimiter and end reception of a certain transmission.

On the other hand, delimeters are usefull when sending data that will always be in the same format every time, but of varying length. This removes the need for certain information to be sent ahead of time, such as payload type and length, and any other overhead the payload may contain. This could reduce the size of the transmission, making for a faster exchange.

\bibliographystyle{plain}
\bibliography{Biblio}

\end{document}

